\subsection{应用：定义在紧空间的乘积上的连续实值函数的逼近}

定理 4. 设 \((\mathbf{X}_t)_{t\in \mathbf{I}}\) 是紧空间的一个族，并令

\[
\mathrm{X} = \prod_ {t \in \mathrm{I}} \mathrm{X} _ {t}.
\]

则 X 上的每个连续实值函数都可用有限多个下列形式的函数之和一致逼近：

\[
(x _ {t}) \rightarrow \prod_ {\alpha \in \mathbf {J}} u _ {\alpha} (x _ {\alpha}),
\]

其中 J 是 I 的（任意）有限子集，而 \(u_{\alpha}\) 是 \(X_{\alpha}\) 上的连续实值函数（对每个 \(\alpha \in J\) ）。

考察 H，即在 X 上连续的“单元函数” \((x_{t}) \rightarrow u_{\alpha}(x_{\alpha})\) (任意 \(\alpha \in I\) )所成的集。这个集分离 X 的点，因为若 \(x = (x_{t})\) 与 \(y = (y_{t})\) 是 X 的任意两个不同点，则存在 \(\alpha \in I\) 使 \(x_{\alpha} \neq y_{\alpha}\) ，并且存在连续实值函数 \(h_{\alpha}\) ，它定义在 \(X_{\alpha}\) 上且使 \(h_{\alpha}(x_{\alpha}) \neq h_{\alpha}(y_{\alpha})\) 。于是函数 \(x \rightarrow h_{\alpha}(\operatorname{pr}_{\alpha} x)\) 属于 H 且在 x 与 y 处取不同的值。由于 H 中函数的每个多项式都具有定理中所述的形式，结论由定理 3 得出。

若诸 \(X_{t}\) 并非全为紧的，定理 4 的结论未必成立（参见习题 9）。

